% Kato，原书第78--105页；承接前一分块。
\origpage{83}{78}
$\CH^r(X)$ 表示余维为 $r$ 的闭链的\term{Chow 群}{Chow group}{groupe de Chow}，而下标 $\mathrm{hom}\sim0$ 表示与零\term{同调等价}{homologically equivalent}{homologiquement équivalent}的元素所成的子空间。

固定素数 $p$，令\translatornote{6}
\[
 V=H^m_{\et}(X\otimes_{\mathbb Q}\overline{\mathbb Q},\mathbb Q_p).
\]
对包含 $p$ 以及 $X$ 的所有坏约化素数的有限素数集 $S$，令
\[
 H^i_S=H^i_{\et}(\mathbb Z[1/S],V).
\]
这里把 $V$ 看作 $\Spec(\mathbb Z[1/S])_{\et}$ 上的光滑 $\mathbb Q_p$ 层（[De$_4$] 1.1.1 b）；对于交换环 $R$，用 $H^i_{\et}(R,\ )$ 表示 $H^i_{\et}(\Spec(R),\ )$。已知 $H^i_S$ 是有限维 $\mathbb Q_p$ 向量空间。定义
\[
 H^i_{\mathrm{lim}}=\varinjlim_S H^i_S.
\]

\paragraph{2.2.5.} 设 $X,m,r,p$ 如 2.2.4，并假设
\[
 \text{(i)}\qquad r>\inf(m,\dim X).
\]
有两个典范映射
\[
 a:\Psi\otimes\mathbb R\longrightarrow\Sigma\otimes\mathbb R,
 \qquad b:\Psi\otimes\mathbb Q_p\longrightarrow H^1_{\mathrm{lim}},
\]
其中 $a$ 是\term{调节子映射}{regulator map}{application régulateur}（[Be$_1$]），$b$ 是\term{陈类映射}{Chern class map}{application classe de Chern}（[So$_1$]）。若不作假设 (i)，则调节子映射的目标群是 $\Sigma\otimes\mathbb R$ 的某个商群，参见 [Be$_1$]。

例如，当 $X=\Spec(K)$，其中 $K$ 为数域，且 $(m,r)=(0,1)$ 时，
\[
 \Sigma=\bigoplus_v\mathbb Q,
 \qquad \Psi=K^\times\otimes\mathbb Q,
 \qquad H^i_S=H^i_{\et}(\mathcal O_K[1/S],\mathbb Q_p(1)),
\]
其中 $v$ 遍历 $K$ 的所有阿基米德位。映射 $a:\Psi\otimes\mathbb R\to\Sigma\otimes\mathbb R$ 是调节子映射
\[
 x\longmapsto(\log\lVert x\rVert_v)_v.
\]
\origpage{84}{79}
映射 $b:\Psi\otimes\mathbb Q_p\to H^1_{\mathrm{lim}}$ 则由 $\Spec(\mathcal O_K[1/S])_{\et}$ 上的 Kummer 正合列
\[
 0\longrightarrow\mathbb Z/p^n(1)\longrightarrow\mathbb G_m
 \xrightarrow{p^n}\mathbb G_m\longrightarrow0
\]
的连接同态
\[
 (\mathcal O_K[1/S])^\times\longrightarrow
 H^1_{\et}(\mathcal O_K[1/S],\mathbb Z/p^n(1))
\]
诱导。

\paragraph{2.2.6.} 设 $X,m,r,p$ 如 2.2.4，$S$ 是包含 $p$ 以及 $X$ 的所有坏约化素数的有限素数集。令 $V$ 如 2.2.4，并令
\[
 V^*=\Hom_{\mathbb Q_p}(V,\mathbb Q_p),
\]
其上的 $\Gal(\overline{\mathbb Q}/\mathbb Q)$ 作用为 $(\sigma,h)\mapsto h\circ\sigma^{-1}$。在猜想 2.2.7 的假设下，下述猜想 2.2.7(3)（分别地，(4)）就是 $L_S(V^*(1),s)$ 在 $s=0$ 处的阶段 1（分别地，阶段 3）。

先对 2.2.7 作几点说明。
\begin{enumerate}[label=(\arabic*)]
\item 猜想 2.2.7(1)--(3) 是 Beilinson 猜想（[Ra]）的一部分；其中涉及 $H^i_{\mathrm{lim}}$ 的部分归于 Jannsen [Ja]，另参见 [So$_3$]。猜想 2.2.7(4) 是 [FP$_1$][FP$_2$] 中一个猜想的一部分。
\item 例如，若 $X$ 在每个素数处都有潜在好约化，则在其余假设 (i)、(ii) 成立时，2.2.7 的假设 (iii) 也成立。
\item 若 $X$ 纯维数为 $d$，由\term{Poincaré 对偶}{Poincaré duality}{dualité de Poincaré}有
\[
 V^*(1)\simeq H^{2d-m}(X\otimes_{\mathbb Q}\overline{\mathbb Q},\mathbb Q_p)(d+1-r).
\]
若 $X$ 还是射影的，由强 Lefschetz 定理（[De$_4$] 4.1）有
\[
 V^*(1)\simeq H^m(X\otimes_{\mathbb Q}\overline{\mathbb Q},\mathbb Q_p)(m+1-r).
\]
\item 在 2.2.7 的假设 (i)、(ii) 下，可以证明\translatornoteref{7}
\[
 H^1_S\xrightarrow{\sim}H^1_{\mathrm{lim}},\qquad H^0_S=0.
\]
\end{enumerate}

\begin{statement}{猜想 2.2.7}
设 $X,m,r,p,S$ 如 2.2.5，并假设下列 (i)--(iii)。令
$V=H^m_{\et}(X\otimes_{\mathbb Q}\overline{\mathbb Q},\mathbb Q_p)(r)$。
\begin{enumerate}[label=(\roman*)]
\item $r>\inf(m,\dim X)$。
\item $(m,r)\ne(1,0),(2\dim X,\dim X+1)$。\translatornote{7}
\item 对所有 $\ell\in S$，有 $P_{\Lambda,\ell}(V^*(1),1)\ne0$。参见 2.2.3(3)、2.2.6(2)。
\end{enumerate}
\origpage{85}{80}
那么：
\begin{enumerate}[label=(\arabic*)]
\item 有同构
\[
 a:\Psi\otimes\mathbb R\xrightarrow{\sim}\Sigma\otimes\mathbb R,
 \qquad b:\Psi\otimes\mathbb Q_p\xrightarrow{\sim}H^1_{\mathrm{lim}},
\]
且 $i\ne1$ 时，$H^i_S=0$。
\item 假设 $L_S(V^*(1),s)$ 可解析延拓到整个 $\mathbb C$。令
\[
 e=\operatorname{ord}_{s=0}L_S(V^*(1),s).
\]
则 $e=\dim_{\mathbb Q}\Sigma$。
\item 在由 $a$ 诱导的同构
\[
 \det_{\mathbb Q}(\Psi)\otimes\mathbb R
 \xrightarrow{\sim}\det_{\mathbb Q}(\Sigma)\otimes\mathbb R
\]
下，$\det_{\mathbb Q}(\Psi)$ 的像恰为
\[
 \left(\lim_{s\to0}s^{-e}L_S(V^*(1),s)\right)\det_{\mathbb Q}(\Sigma).
\]
\item 假设 $p\ne2$。令
\[
 \Sigma_{\mathbb Z}\overset{\mathrm{def}}=
 \bigl(H^m(X(\mathbb C),\mathbb Z)\otimes\mathbb Z(2\pi i)^{r-1}\bigr)^+,
\]
并定义有限生成 $\mathbb Z_p$ 模
\[
 H^i_{S,\mathbb Z_p}\overset{\mathrm{def}}=
 H^i_{\et}(\mathbb Z[1/S],T),\qquad
 T=H^m_{\et}(X\otimes_{\mathbb Q}\overline{\mathbb Q},\mathbb Z_p)(r).
\]
取 $\det_{\mathbb Z}(\Sigma_{\mathbb Z})$ 的一个 $\mathbb Z$ 基 $\gamma$。令 $\omega_\gamma$ 为 $\det_{\mathbb Q}(\Psi)$ 中在 (3) 的同构下映为
\[
 \left(\lim_{s\to0}s^{-e}L_S(V^*(1),s)\right)\gamma
\]
的元素；由 (3)，此元素存在。再令 $z_\gamma$ 是 $\omega_\gamma$ 在 $b$ 诱导的同构
\[
 \det_{\mathbb Q}(\Psi)\otimes\mathbb Q_p
 \xrightarrow{\sim}\det_{\mathbb Q_p}(H^1_{\mathrm{lim}})
\]
下的像。那么
\[
 [\det_{\mathbb Z_p}(H^1_{S,\mathbb Z_p}):\mathbb Z_p z_\gamma]
 =\#H^0_{S,\mathbb Z_p}\cdot\#H^2_{S,\mathbb Z_p}.
\]
换言之（2.1.5），通过同构
\[
 \{\det_{\mathbb Z_p}R\Gamma(\mathbb Z[1/S],T)\}^{-1}
 \otimes_{\mathbb Z_p}\mathbb Q_p
 \simeq\det_{\mathbb Q_p}(H^1_{\mathrm{lim}}),
\]
$z_\gamma$ 是 $\{\det_{\mathbb Z_p}R\Gamma(\mathbb Z[1/S],T)\}^{-1}$ 的一个 $\mathbb Z_p$ 基。
\end{enumerate}
\end{statement}

\origpage{86}{81}
\paragraph{2.2.8.} 关于猜想 2.2.7(3)，即阶段 1，的一些已知结果，见 [Ra]。

这里回顾 $X=\Spec(K)$（$K$ 为数域）、$m=0$、$r\ge2$ 的情形。此时 2.2.7 的 (i)、(ii)、(iii) 都成立，并且 (1)--(3) 中的所有结论均已得到证明。事实上，在 (1) 中，$a$ 的双射性由 Borel [Bo] 证明，$b$ 的双射性由 Soulé [So$_2$] 证明。注意
\[
 L_S(V^*(1),s)=\zeta_K^{(S)}(s+1-r),
\]
其中 $\zeta_K^{(S)}(s)$ 是去掉 $S$ 上方素理想处 Euler 因子的 $K$ 的 Dedekind $\zeta$ 函数。因此 (2) 是经典结果，(3) 则是 Borel [Bo] 的结果。已知（[Be$_1$]），若 $q,i\in\mathbb Z$，$q\ne2i-1$ 且 $(q,i)\ne(0,0)$，则 $(K_q(K)\otimes\mathbb Q)^{(i)}=0$。

\paragraph{2.2.9.} 下面回顾 $X=\Spec(K)$（$K$ 为数域）、$m=0$、$r\ge2$ 时猜想 2.2.7(4)，即阶段 3，的情况。

先设 $K$ 为全实域，且 $r$ 为偶数。此时 $\Sigma,\Psi$ 以及所有 $i\in\mathbb Z$ 对应的 $H^i_S$ 均为零。由 2.1.5，2.2.7(4) 就是 Lichtenbaum 猜想\translatornote{8}
\[
 \zeta_K(1-r)_{(p)}
 =\#H^2(\mathcal O_K[1/S],\mathbb Z_p(r))
 \cdot\#H^1(\mathcal O_K[1/S],\mathbb Z_p(r))^{-1}.
\]
对某些特殊的 $K$，Mazur--Wiles [MW] 已证明这一猜想；一般情形由 Wiles [Wi$_2$] 证明。

对一般的 $r\ge2$，当 $K=\mathbb Q$ 时，2.2.7(4) 也成立。[BK] §6 假设了 [BK] 6.2 的一个猜想，利用 Mazur--Wiles 对岩泽主猜想的证明推出这一结论。这个附加猜想后来由 Beilinson [Be] 证明，即 3.4.9 末尾将介绍的 Beilinson 的结果。看来，当 $K/\mathbb Q$ 为阿贝尔扩张时，也可以用类似的方法证明。

\origpage{87}{82}
\section*{§2.3. Birch--Swinnerton-Dyer 猜想}
\addcontentsline{toc}{section}{§2.3. Birch--Swinnerton-Dyer 猜想}
在阿贝尔簇的 $L$ 函数于 $s=1$ 处不为零的假设下，我们给出 Fontaine 与 Perrin-Riou [FP$_1$][FP$_2$] 对 Birch--Swinnerton-Dyer 猜想的重新表述。这里不考虑 $2$ 准素部分。

\paragraph{2.3.1.} 一般地，设 $X$ 为 $\mathbb Q$ 上的固有光滑概形，固定整数 $m,r\ge0$ 以及素数 $p$。令 $S$ 为包含 $p$ 和 $X$ 的所有坏约化素数的有限素数集。如 §2.2，令
\[
 \Sigma=H^m(X(\mathbb C),\mathbb Q(2\pi i)^{r-1}),\qquad
 V=H^m_{\et}(X\otimes_{\mathbb Q}\overline{\mathbb Q},\mathbb Q_p)(r),
 \qquad H^i_S=H^i(\mathbb Z[1/S],V).
\]
上式 $\Sigma$ 的定义未印上标 $+$。\footnote{译注：原书如此；比较 2.2.4 的定义及下文周期映射中的不动部分。}

令
\[
 \Omega=\operatorname{fil}^rH^m_{\dR}(X/\mathbb Q),
\]
其中 $\operatorname{fil}^\bullet$ 为\term{Hodge 滤过}{Hodge filtration}{filtration de Hodge}，定义为
\[
 \operatorname{fil}^iH^m_{\dR}(X/\mathbb Q)
 =\im\bigl(H^m(X,\Omega^{\ge i}_{X/\mathbb Q})\longrightarrow H^m_{\dR}(X/\mathbb Q)\bigr)
 \quad(i\in\mathbb Z),
\]
这里 $\Omega^{\ge i}_{X/\mathbb Q}$ 是 de Rham 复形 $\Omega^\bullet_{X/\mathbb Q}$ 的次数 $\ge i$ 的部分。

于是有典范同态
\[
 a:\Omega\otimes\mathbb R\longrightarrow\Sigma\otimes\mathbb R,
 \qquad b:H^1_S\longrightarrow\Omega\otimes\mathbb Q_p,
\]
它们分别由\term{周期积分}{period integral}{intégrale de période}理论和 $p$ 进周期理论定义。注意，相对于 §2.2 中的 $a,b$，这里 $\Omega$ 扮演彼处 $\Psi$ 的角色，但 $b$ 的方向相反。映射 $a$ 由周期同构
\[
 H^m_{\dR}(X/\mathbb Q)\otimes\mathbb C
 \xrightarrow{\sim}H^m(X(\mathbb C),\mathbb Q)\otimes\mathbb C
\]
（参见 1.2.2）按如下方式诱导：
\begin{align*}
 \Omega\otimes\mathbb R&\hookrightarrow H^m_{\dR}(X/\mathbb Q)\otimes\mathbb R
 \simeq\bigl(H^m(X(\mathbb C),\mathbb Q)\otimes\mathbb C\bigr)^+\\
 &\longrightarrow
 \frac{\bigl(H^m(X(\mathbb C),\mathbb Q)\otimes\mathbb C\bigr)^+}
 {\bigl(H^m(X(\mathbb C),\mathbb Q)\otimes\mathbb R(2\pi i)^r\bigr)^+}
 \longleftarrow\Sigma\otimes\mathbb R.
\end{align*}

\origpage{88}{83}
映射 $b$ 的定义稍复杂。它是下列映射的复合：
\begin{align*}
 H^1(\mathbb Z[1/S],V)&\longrightarrow H^1(\mathbb Q_p,V)
 \xrightarrow{\exp^*}\operatorname{fil}^rH^m_{\dR}((X\otimes_{\mathbb Q}\mathbb Q_p)/\mathbb Q_p)\\
 &\hspace{7em}\simeq\operatorname{fil}^rH^m_{\dR}(X/\mathbb Q)\otimes\mathbb Q_p.
\end{align*}
这里的同态 $\exp^*$ 称为\term{对偶指数映射}{dual exponential map}{application exponentielle duale}，将在第 II 章 §1.2 中用 $p$ 进周期理论定义。

\paragraph{2.3.2.} 现在设 $X$ 是 $\mathbb Q$ 上的阿贝尔簇，取 $m=r=1$。那么 $\Omega=\Gamma(X,\Omega^1_{X/\mathbb Q})$，此时 $a$ 为双射。在这个特殊情形中，上面的 $\exp^*$ 无须借助 $p$ 进周期理论就能定义，详见后面的 2.3.6。

由于有非典范同构 $V^*(1)\simeq V$，
\[
 L_S(V^*(1),s)=L_S(V,s)=L_S(H^1(X),s+1).
\]
假设
\begin{equation*}
 \tag{*}
 \begin{gathered}
 L_S(H^1(X),s)\text{ 可作为全纯函数解析延拓到整个 }\mathbb C,\\
 L_S(H^1(X),1)\ne0.
 \end{gathered}
\end{equation*}
那么，关于 $L_S(H^1(X),1)$ 的阶段 1 是下面的 2.3.3(1)，阶段 3 是 (2) 和 (3)。

\begin{statement}{猜想 2.3.3（$L_S(H^1(X),1)\ne0$ 时的 Birch--Swinnerton-Dyer 猜想）}
设 $X$ 是 $\mathbb Q$ 上的阿贝尔簇，并假设 (*)。
\begin{enumerate}[label=(\arabic*)]
\item 在由 $a$ 诱导的同构
\[
 \det_{\mathbb Q}(\Omega)\otimes\mathbb R
 \xrightarrow{\sim}\det_{\mathbb Q}(\Sigma)\otimes\mathbb R
\]
下，$\det_{\mathbb Q}(\Omega)$ 的像恰为 $L_S(H^1(X),1)\det_{\mathbb Q}(\Sigma)$。
\item $b$ 为同构，且 $i\ne1$ 时 $H^i_S=0$。
\item 假设 $p\ne2$。令 $\Sigma_{\mathbb Z}=H^1(X(\mathbb C),\mathbb Z)^+$，并定义有限生成 $\mathbb Z_p$ 模
\[
 H^i_{S,\mathbb Z_p}=H^i_{\et}(\mathbb Z[1/S],T),\qquad
 T=H^1_{\et}(X\otimes_{\mathbb Q}\overline{\mathbb Q},\mathbb Z_p)(1).
\]
\origpage{89}{84}
取 $\det_{\mathbb Z}(\Sigma_{\mathbb Z})$ 的一个 $\mathbb Z$ 基 $\gamma$，令 $\omega_\gamma\in\det_{\mathbb Q}(\Omega)$ 为 (1) 中的同构下映为 $L_S(H^1(X),1)\gamma$ 的元素；由 (1)，它存在。令 $z_\gamma\in\det_{\mathbb Q_p}(H^1_S)$ 为由 $b$ 诱导的同构
\[
 \det_{\mathbb Q_p}(H^1_S)\xrightarrow{\sim}\det_{\mathbb Q}(\Omega)\otimes\mathbb Q_p
\]
下映为 $\omega_\gamma$ 的元素。那么
\[
 [\det_{\mathbb Z_p}(H^1_{S,\mathbb Z_p}):\mathbb Z_pz_\gamma]
 =\#H^2_{S,\mathbb Z_p}.
\]
换言之，$z_\gamma$ 是 $\{\det_{\mathbb Z_p}R\Gamma(\mathbb Z[1/S],T)\}^{-1}$ 的一个 $\mathbb Z_p$ 基。
\end{enumerate}
\end{statement}

\paragraph{备注 2.3.4.} (1) Birch--Swinnerton-Dyer 猜想通常以另一种形式表述（[Ta$_1$]）。

(2) 关于椭圆曲线的 Birch--Swinnerton-Dyer 猜想的结果，参见 Kolyvagin [Ko] 与 Rubin [Ru]。

\paragraph{2.3.5.} 在猜想 2.2.7 与 2.3.3 中，我们得到可逆 $\mathbb Z_p$ 模
\[
 \{\det_{\mathbb Z_p}R\Gamma(\mathbb Z[1/S],T)\}^{-1}
 \otimes_{\mathbb Z}\{\det_{\mathbb Z}(\Sigma_{\mathbb Z})\}^{-1}
\]
的一个猜想中的典范 $\Lambda$ 基 $z_\gamma\otimes\gamma^{-1}$。\footnote{译注：原书如此，称为 $\Lambda$ 基；本段所述模的系数环为 $\mathbb Z_p$。}在相差一个符号的意义下，这个基与我们的岩泽主猜想 3.2.2 中（猜想存在的）$\mathbb Z_p$ 层 $T$ 的“\term{zeta 元素}{zeta element}{élément zêta}”一致。

\paragraph{2.3.6.} 设 $X$ 为 $\mathbb Q_p$ 上的阿贝尔簇，
\[
 V=H^1(X\otimes_{\mathbb Q_p}\overline{\mathbb Q}_p,\mathbb Q_p)(1).
\]
映射 $\exp^*:H^1(\mathbb Q_p,V)\to\Gamma(X,\Omega^1_{X/\mathbb Q_p})$ 的定义如下。注意，$\Gamma(X,\Omega^1_{X/\mathbb Q_p})$ 可识别为 $X$ 在原点处的\term{切空间}{tangent space}{espace tangent} $\operatorname{Lie}(X)$ 的 $\mathbb Q_p$ 对偶。$p$ 进 Lie 群 $X(\mathbb Q_p)$ 有指数映射
\origpage{90}{85}
\[
 \exp:\operatorname{Lie}(X)\xrightarrow{\sim}X(\mathbb Q_p)\otimes\mathbb Q.
\]
另一方面，定义 ${}_{p^n}X$ 的正合列
\[
 0\longrightarrow{}_{p^n}X\longrightarrow X\xrightarrow{p^n}X\longrightarrow0
\]
给出 Galois 上同调的连接映射
\[
 X(\mathbb Q_p)\longrightarrow H^1(\mathbb Q_p,{}_{p^n}X).
\]
取 $\mathbb Q\otimes\varprojlim$，得到
\[
 X(\mathbb Q_p)\longrightarrow H^1(\mathbb Q_p,V_pX),\qquad
 V_pX=T_pX\otimes_{\mathbb Z_p}\mathbb Q_p,
\]
其中 $T_pX$ 为 $X$ 的 $p$ 进\term{Tate 模}{Tate module}{module de Tate}。仍用 $\exp$ 表示此映射与上述指数映射的复合：
\[
 \exp:\operatorname{Lie}(X)\longrightarrow H^1(\mathbb Q_p,V_pX).
\]
由局部 Tate 对偶（参见第 II 章 1.4），以及典范同构 $V_pX\simeq V^*(1)$，$H^1(\mathbb Q_p,V)$ 典范同构于 $H^1(\mathbb Q_p,V_pX)$ 的 $\mathbb Q_p$ 对偶。因此，对最后的 $\exp$ 取 $\mathbb Q_p$ 对偶，即得
\[
 \exp^*:H^1(\mathbb Q_p,V)\longrightarrow\Gamma(X,\Omega^1_{X/\mathbb Q_p}).
\]

\section*{§3. 广义岩泽主猜想}
\addcontentsline{toc}{section}{§3. 广义岩泽主猜想}
我们将表述广义\term{岩泽主猜想}{Iwasawa main conjecture}{conjecture principale d'Iwasawa} 3.2.2。在 §3.3 中讨论它与经典岩泽主猜想的关系，在 §3.4 中讨论由此产生的重要问题；§3.5 则讨论有限域上的簇的情形。

在本节中，$p$ 表示一个固定的素数。

\origpage{91}{86}
\subsection*{§3.1. 预备知识}
\paragraph{3.1.1.} 我们考虑满足下列两种条件之一的交换\term{拓扑环}{topological ring}{anneau topologique} $\Lambda$：
\begin{enumerate}[label=(\arabic*)]
\item $\Lambda$ 是 pro-$p$ 环，即一个紧环，并且是阶为 $p$ 的幂的有限环的逆极限；
\item $\Lambda$ 是 $\mathbb Q_p$ 的有限扩张的有限直积，因而是一个 $\mathbb Q_p$ 代数。
\end{enumerate}

\paragraph{3.1.2.} 设 $\Lambda$ 如 3.1.1。如下定义范畴 $D_{\mathrm{ctf}}(\Lambda)$，以及任意 Noether 概形 $X$ 对应的范畴 $D_{\mathrm{ctf}}(X,\Lambda)$。

(1) 先设 $\Lambda$ 为有限环。令 $D_{\mathrm{ctf}}(\Lambda)$ 为 $\Lambda$ 模的\term{导出范畴}{derived category}{catégorie dérivée} $D(\Lambda)$ 中由\term{完美复形}{perfect complex}{complexe parfait}（2.1.1）组成的满子范畴。若 $X$ 为 Noether 概形，令 $D(X,\Lambda)$ 为 $X_{\et}$ 上的 $\Lambda$ 模层的导出范畴，$D_{\mathrm{ctf}}(X,\Lambda)$ 为其中满足下列 (i)、(ii) 的对象 $\mathcal F$ 所成的满子范畴：
\begin{enumerate}[label=(\roman*)]
\item 对所有 $q\in\mathbb Z$，上同调层 $\mathcal H^q(\mathcal F)$ 都是\term{可构造层}{constructible sheaf}{faisceau constructible}（[SGA 4]），且除有限多个 $q$ 外均为零；
\item 对任意 $x\in X$，茎 $\mathcal F_{\bar x}$ 是 $\Lambda$ 上的完美复形。
\end{enumerate}

(2) 若 $\Lambda$ 满足 3.1.1(1)，则如下定义 $D_{\mathrm{ctf}}(\Lambda)$，以及对 Noether 概形 $X$ 的 $D_{\mathrm{ctf}}(X,\Lambda)$。其对象是一个系统 $(\mathcal F_I,\rho_{I,J})$，其中 $I$ 遍历 $\Lambda$ 的开理想，$\mathcal F_I$ 为 $D_{\mathrm{ctf}}(\Lambda/I)$（分别地，$D_{\mathrm{ctf}}(X,\Lambda/I)$）的对象。对每一对满足 $I\supset J$ 的开理想，给定同构
\[
 \rho_{I,J}:\mathcal F_J\otimes^{\mathbf L}_{\Lambda/J}\Lambda/I
 \xrightarrow{\sim}\mathcal F_I,
\]
并要求 $\rho_{I,I}=\id$，以及当 $I\supset I'\supset I''$ 时
\[
 \rho_{I,I'}\circ\rho_{I',I''}=\rho_{I,I''}.
\]
这里 $\otimes^{\mathbf L}$ 表示导出范畴中的张量积。
\origpage{92}{87}
$D_{\mathrm{ctf}}(\Lambda)$（分别地，$D_{\mathrm{ctf}}(X,\Lambda)$）中的态射以自然方式定义。

(3) 最后，设 $\Lambda$ 满足 3.1.1(2)。写成 $F=\prod_iF_i$，其中每个 $F_i$ 为 $\mathbb Q_p$ 的有限扩张；令 $\mathcal O_i$ 为 $F_i$ 的整数环，$\mathcal O_\Lambda=\prod_i\mathcal O_i$。定义
\[
 D_{\mathrm{ctf}}(\Lambda)=D_{\mathrm{ctf}}(\mathcal O_\Lambda)\otimes\mathbb Q,
 \qquad
 D_{\mathrm{ctf}}(X,\Lambda)=D_{\mathrm{ctf}}(X,\mathcal O_\Lambda)\otimes\mathbb Q.
\]
这里，对于加性范畴 $\mathcal C$，范畴 $\mathcal C\otimes\mathbb Q$ 与 $\mathcal C$ 有相同的对象，而其 Hom 群为 $\mathcal C$ 的相应 Hom 群与 $\mathbb Q$ 的张量积。

下面的结果是已知的。
\begin{statement}{3.1.3}
设 $\Lambda$ 为阶为 $p$ 的幂的有限环。

(1) 设 $Z$ 为维数 $\le1$ 的正则 Noether 概形，且 $p$ 在 $Z$ 上可逆；设 $f:X\to Y$ 为 $Z$ 上有限型概形间的 $Z$ 态射。那么 $Rf_*$ 把 $D_{\mathrm{ctf}}(X,\Lambda)$ 映入 $D_{\mathrm{ctf}}(Y,\Lambda)$。

(2) 考虑下列情形：
\begin{enumerate}[label=(\roman*)]
\item $X$ 为 $\mathbb Z[1/p]$ 上的有限型概形，且 $p\ne2$；
\item $X$ 为 $\mathbb R$ 上的有限型概形，且 $p\ne2$；
\item $X$ 为 $\mathbb Q_\ell$ 上的有限型概形，其中 $\ell$ 为一个素数。
\end{enumerate}
设 $\mathcal F\in D_{\mathrm{ctf}}(X,\Lambda)$。在 (i)--(iii) 的各情形中，$R\Gamma(X,\mathcal F)\in D_{\mathrm{ctf}}(\Lambda)$。而且，对于任意阶为 $p$ 的幂的有限环 $\Lambda'$ 以及任意同态 $\Lambda\to\Lambda'$，有
\[
 R\Gamma(X,\mathcal F)\otimes^{\mathbf L}_\Lambda\Lambda'
 \xrightarrow{\sim}R\Gamma(X,\mathcal F'),
 \qquad\mathcal F'=\mathcal F\otimes^{\mathbf L}_\Lambda\Lambda'.
\]
在情形 (i) 中，
\begin{equation*}\tag{*}
 \rank_\Lambda R\Gamma(X,\mathcal F)
 =-\rank_\Lambda R\Gamma(X\otimes\mathbb R,\mathcal F(-1)).
\end{equation*}
在情形 (iii) 中，若 $\ell=p$，则
\begin{equation*}\tag{**}
 \rank_\Lambda R\Gamma(X,\mathcal F)
 =-\rank_\Lambda R\Gamma(X\otimes_{\mathbb Q_p}\overline{\mathbb Q}_p,\mathcal F);
\end{equation*}
若 $\ell\ne p$，则 $\rank_\Lambda R\Gamma(X,\mathcal F)=0$。
\end{statement}

\origpage{93}{88}
\begin{proof}
(1) 见 Deligne [De$_2$] 和 [SGA 4] 第 17 章。(2) 中除秩以外的部分，可用同样的方法证明：利用当 $X=\Spec(\mathbb Z[1/p])$、$\Spec(\mathbb R)$ 或 $\Spec(\mathbb Q_\ell)$ 时，所有 $H^i(X,\mathcal F)$ 都有限，且除有限多个 $i$ 外均为零这一事实（[Ma]、[Se]）。关于秩的问题，可归结到 $\Lambda$ 为域的情形，再由 Tate [Ta$_1$] 定理 2.2 与 Serre [Se] 第 II 章 §5.7 得到。
\end{proof}

\paragraph{3.1.4.} 在 3.1.3 中，将“$\Lambda$ 为阶为 $p$ 的幂的有限环”替换为“$\Lambda$ 如 3.1.1(1)”。对于 $\mathcal F=(\mathcal F_I)_I\in D_{\mathrm{ctf}}(X,\Lambda)$，在 3.1.3(1)（分别地，(2)）中，定义
\[
 Rf_*(\mathcal F)=(Rf_*(\mathcal F_I))_I
 \quad\text{（分别地， }R\Gamma(X,\mathcal F)=(R\Gamma(X,\mathcal F_I))_I\text{）},
\]
它是 $D_{\mathrm{ctf}}(Y,\Lambda)$（分别地，$D_{\mathrm{ctf}}(\Lambda)$）的对象。

若将“$\Lambda$ 为阶为 $p$ 的幂的有限环”替换为“$\Lambda$ 如 3.1.1(2)”，则对 $\mathcal F\in D_{\mathrm{ctf}}(X,\mathcal O_\Lambda)$，在 3.1.3(1)（分别地，(2)）中定义
\[
 Rf_*(\mathcal F\otimes\mathbb Q)=Rf_*(\mathcal F)\otimes\mathbb Q
 \quad\text{（分别地， }R\Gamma(X,\mathcal F\otimes\mathbb Q)=R\Gamma(X,\mathcal F)\otimes\mathbb Q\text{）}.
\]

\subsection*{§3.2. 广义岩泽主猜想}
令 $p$ 为一个固定素数。

\paragraph{3.2.1.} 我们考虑满足下列条件的三元组 $(X,\Lambda,\mathcal F)$：
\begin{enumerate}[label=(\roman*)]
\item $X$ 为 $\mathbb Z[1/p]$ 上的有限型概形；
\item $\Lambda$ 如 3.1.1；
\item $\mathcal F$ 为 $D_{\mathrm{ctf}}(X,\Lambda)$ 的对象。
\end{enumerate}
对这样的三元组，定义\term{可逆模}{invertible module}{module inversible}
\[
 \Delta_\Lambda(X,\mathcal F)
 =\{\det_\Lambda R\Gamma(X,\mathcal F)\}^{-1}
 \otimes_\Lambda
 \{\det_\Lambda R\Gamma(X\otimes_{\mathbb Z}\mathbb R,\mathcal F(-1))\}^{-1}.
\]

\origpage{94}{89}
\begin{statement}{猜想 3.2.2}
设 $p$ 为奇素数。那么，对任意如 3.2.1 的三元组 $(X,\Lambda,\mathcal F)$，可以定义 $\Delta_\Lambda(X,\mathcal F)$ 的一个 $\Lambda$ 基 $z_\Lambda(X,\mathcal F)$，称为 $(X,\Lambda,\mathcal F)$ 的 zeta 元素，使之满足下列 (i)--(v)。条件 (v) 涉及 zeta 函数值和“动机”一词，这里仅作部分陈述；在第 III 章 §2 介绍动机的预备知识后，将详细给出该条件。
\begin{enumerate}[label=(\roman*)]
\item 设 $(X,\Lambda,\mathcal F)$、$(X,\Lambda',\mathcal F')$ 是具有相同 $X$ 的三元组，均满足 3.2.1。设有同态 $\Lambda\to\Lambda'$，以及 $\Lambda'$ 同构
\[
 h:\mathcal F\otimes^{\mathbf L}_\Lambda\Lambda'\xrightarrow{\sim}\mathcal F'.
\]
这里允许 $\Lambda$ 满足 3.1.1(1)，而 $\Lambda'$ 满足 3.1.1(2)。那么，由 $h$ 诱导的同构（3.1.3(2)）
\[
 \Delta_\Lambda(X,\mathcal F)\otimes^{\mathbf L}_\Lambda\Lambda'
 \xrightarrow{\sim}\Delta_{\Lambda'}(X,\mathcal F')
\]
把 $z_\Lambda(X,\mathcal F)\otimes1$ 映为 $z_{\Lambda'}(X,\mathcal F')$。
\item 设 $(X,\Lambda,\mathcal F)$ 如 3.2.1，$Y$ 为 $\mathbb Z[1/p]$ 上的有限型概形，$f:X\to Y$ 为态射。那么，在 $\Delta_\Lambda(X,\mathcal F)=\Delta_\Lambda(Y,Rf_*\mathcal F)$ 中，
\[
 z_\Lambda(X,\mathcal F)=z_\Lambda(Y,Rf_*\mathcal F).
\]
\item 假设 $\Lambda$ 为有限环。设 $(X,\Lambda,\mathcal F)$、$(X,\Lambda,\mathcal F')$、$(X,\Lambda,\mathcal F'')$ 均如 3.2.1，具有相同的 $X,\Lambda$。给定内射 $\Lambda$ 层复形的正合列
\begin{equation*}\tag{*}
 0\longrightarrow I'^\bullet\longrightarrow I^\bullet\longrightarrow I''^\bullet\longrightarrow0,
\end{equation*}
其中 $I^\bullet,I'^\bullet,I''^\bullet$ 分别代表 $\mathcal F,\mathcal F',\mathcal F''$。那么，通过同构
\[
 \Delta_\Lambda(X,\mathcal F)\simeq
 \Delta_\Lambda(X,\mathcal F')\otimes^{\mathbf L}_\Lambda\Delta_\Lambda(X,\mathcal F''),
\]
有 $z_\Lambda(X,\mathcal F)=z_\Lambda(X,\mathcal F')\otimes z_\Lambda(X,\mathcal F'')$。这里的同构由下列复形正合列诱导：
\[
 0\to\Gamma(X,I'^\bullet)\to\Gamma(X,I^\bullet)\to\Gamma(X,I''^\bullet)\to0,
\]
\[
 \begin{gathered}
 0\to\Gamma(X\otimes\mathbb R,I'^\bullet(-1))
 \to\Gamma(X\otimes\mathbb R,I^\bullet(-1))\\
 \to\Gamma(X\otimes\mathbb R,I''^\bullet(-1))\to0.
 \end{gathered}
\]
\origpage{95}{90}
\item 若 $\Lambda$ 是有限环，$X$ 为某个素数 $\ell\ne p$ 对应的 $\mathbb F_\ell$ 上的概形，则如下定义 $z_\Lambda(X,\mathcal F)$。令 $\overline X=X\otimes_{\mathbb F_\ell}\overline{\mathbb F}_\ell$，其中 $\overline{\mathbb F}_\ell$ 为 $\mathbb F_\ell$ 的代数闭包。取代表 $\mathcal F$ 的内射 $\Lambda$ 层复形 $I^\bullet$，考虑复形正合列
\begin{equation*}\tag{*}
 0\to\Gamma(X,I^\bullet)\to\Gamma(\overline X,I^\bullet)
 \xrightarrow{1-\sigma_\ell}\Gamma(\overline X,I^\bullet)\to0.
\end{equation*}
这里 $\sigma_\ell$ 为 $\Gal(\overline{\mathbb F}_\ell/\mathbb F_\ell)$ 中的\term{算术 Frobenius}{arithmetic Frobenius}{Frobenius arithmétique}，即对 $x\in\overline{\mathbb F}_\ell$，$\sigma_\ell(x)=x^\ell$。考虑同构
\begin{equation*}\tag{**}
 \Delta_\Lambda(X,\mathcal F)
 \simeq\det_\Lambda R\Gamma(\overline X,\mathcal F)
 \otimes^{\mathbf L}_\Lambda\{\det_\Lambda R\Gamma(\overline X,\mathcal F)\}^{-1}
 \simeq\Lambda.
\end{equation*}
第一个同构由 (*) 给出，第二个由两份 $R\Gamma(\overline X,\mathcal F)$ 间的恒等映射诱导。定义 $z_\Lambda(X,\mathcal F)$ 为 $1\in\Lambda$ 经 (**) 的复合同构所得的像。
\item 这是关于与 $\mathbb Q$ 上动机的 zeta 值之间关系的条件。完整陈述见第 III 章 §2；下文 3.2.6 给出若干特殊情形。
\end{enumerate}
\end{statement}

\paragraph{备注 3.2.3.}
(1) 也可以表述包含 $p=2$ 情形的猜想，参见 3.4.14。

(2) 由条件 (ii)，要找出 3.2.2 中的 zeta 元素，实际上只须在 $X=\Spec(\mathbb Z[1/p])$ 的情形中找出它们。

(3) 粗略地说，3.2.2(v) 是指：若 $X=\Spec(\mathbb Z[1/S])$，其中 $S$ 为包含 $p$ 的有限素数集，而 $\mathcal F$ 是 $X$ 上来自 $\mathbb Q$ 上某个动机的光滑层，则 zeta 元素在相差一个符号的意义下“表达”某个 zeta 值。这个 $\pm1$ 的符号应当确定下来，但我尚未充分理解它。

(4) 粗略地说，条件 3.2.2(v) 来自把 Fontaine 与 Perrin-Riou [FP$_1$][FP$_2$] 理论中的动机推广为“带系数的动机”。

\origpage{96}{91}
\paragraph{3.2.4.} 为了部分说明 (v)，先稍谈 $\mathbb Q$ 上的\term{动机}{motif}{motif}。关于 Grothendieck 提出的“motif”思想，参见 [Ra][FP$_1$] 等；在当前的数学中，它的定义还没有定论。事实上，我们不讨论动机究竟是什么。无论如何，$\mathbb Q$ 上的动机 $M$（更准确地说，是\term{混合动机}{mixed motive}{motif mixte}）是某种“来自 $\mathbb Q$ 上代数几何”的“东西”，它给出带有下列结构的三元组 $(M_B,M_{\dR},M_p)$。以下固定素数 $p$。
\begin{enumerate}[label=(\arabic*)]
\item $M_B$ 称为 $M$ 的\term{Betti 实现}{Betti realization}{réalisation de Betti}，是带有 $\Gal(\mathbb C/\mathbb R)$ 作用的有限维 $\mathbb Q$ 向量空间。
\item $M_{\dR}$ 称为 $M$ 的\term{de Rham 实现}{de Rham realization}{réalisation de de Rham}，是有限维 $\mathbb Q$ 向量空间，带有递减滤过 $(M^i_{\dR})_{i\in\mathbb Z}$，并且 $i\ll0$ 时 $M^i_{\dR}=M_{\dR}$，$i\gg0$ 时 $M^i_{\dR}=0$。
\item $M_p$ 称为 $M$ 的\term{$p$ 进平展实现}{p-adic étale realization}{réalisation étale p-adique}，是 $\Spec(\mathbb Q)_{\et}$ 上的光滑 $\mathbb Q_p$ 层，并且来自某个有限素数集 $S$ 对应的 $\Spec(\mathbb Z[1/S])_{\et}$ 上的光滑 $\mathbb Q_p$ 层。固定 $\mathbb Q$ 的代数闭包 $\overline{\mathbb Q}$ 后，$M_p$ 就是带有 $\Gal(\overline{\mathbb Q}/\mathbb Q)$ 连续作用、且在几乎所有素数处不分歧的有限维 $\mathbb Q_p$ 向量空间；这种解释由取 $\Spec(\overline{\mathbb Q})$ 处的茎得到。
\item $M_B$ 与 $M_{\dR}$ 的关系：给定一个与 $\Gal(\mathbb C/\mathbb R)$ 作用相容的 $\mathbb C$ 向量空间同构 $M_B\otimes\mathbb C\simeq M_{\dR}\otimes\mathbb C$。这里 $\sigma\in\Gal(\mathbb C/\mathbb R)$ 在前者上按 $\sigma\otimes\sigma$ 作用，在后者上按 $1\otimes\sigma$ 作用。
\item $M_B$ 与 $M_p$ 的关系：给定一个与 $\Gal(\mathbb C/\mathbb R)$ 作用相容的 $\mathbb Q_p$ 向量空间同构
\origpage{97}{92}
\[
 M_B\otimes\mathbb Q_p\simeq\Gamma(\Spec(\mathbb C),M_p).
\]
\item $M_{\dR}$ 与 $M_p$ 的关系：给定滤过 $\mathbb Q_p$ 向量空间的同构 $M_{\dR}\otimes\mathbb Q_p\simeq D_{\dR}(M_p)$。这里所用函子 $D_{\dR}$ 的定义，将在第 II 章 §1.1 中回顾。
\end{enumerate}
例如，若 $X$ 为 $\mathbb Q$ 上的光滑固有概形，人们相信对每一对 $m,r\in\mathbb Z$，都存在动机 $M=H^m(X)(r)$，其实现分别为
\begin{align*}
 M_B&=H^m(X(\mathbb C),\mathbb Q)\otimes\mathbb Q(2\pi i)^r,
 \quad\text{带自然的 }\Gal(\mathbb C/\mathbb R)\text{ 作用},\\
 M_{\dR}&=H^m_{\dR}(X/\mathbb Q),\qquad
 M^i_{\dR}=\operatorname{fil}^{i+r}H^m_{\dR}(X/\mathbb Q),\\
 M_p&=H^m_{\et}(X\otimes_{\mathbb Q}\overline{\mathbb Q},\mathbb Q_p)(r),
 \quad\text{带自然的 }\Gal(\overline{\mathbb Q}/\mathbb Q)\text{ 作用}.
\end{align*}
这里 $\operatorname{fil}$ 为 Hodge 滤过（2.3.1）。

\paragraph{3.2.5.} 人们相信，$\mathbb Q$ 上的动机构成一个阿贝尔 $\mathbb Q$ 范畴，并且上述实现函子都是正合且忠实的。所谓 $\mathbb Q$ 范畴，是指 Hom 是 $\mathbb Q$ 向量空间，且态射的复合为 $\mathbb Q$ 双线性。

设 $A$ 为 $\mathbb Q$ 上的交换环，且为 $\mathbb Q$ 的有限扩张的有限直积。所谓“$\mathbb Q$ 上的 $A$ 动机”，是指带有 $A$ 作用的 $\mathbb Q$ 上动机。

人们相信，若 $M$ 为 $A$ 动机，$S$ 为包含 $p$ 的有限素数集，使 $\Gal(\overline{\mathbb Q}/\mathbb Q)$ 在 $M_p$ 上的作用在 $S$ 外不分歧，则 $(A\otimes\mathbb Q_p,M_p,S,A)$ 满足 2.2.3 中的 (i)、(ii)，且 $L_{A,S}(M_p,s)$ 可延拓为整个 $\mathbb C$ 上的亚纯函数。

人们还相信，$\mathbb Q$ 上的 $A$ 动机范畴具有 $A$ 上的张量积、$A$ 对偶 $\Hom_A(\ ,A)$ 和 Tate 扭转等函子，并且这些函子与实现函子交换。

\origpage{98}{93}
以下假定上述信念全部成立。

\paragraph{3.2.6.} 条件 3.2.2(v) 具有如下形式。设给定 $\mathbb Q$ 上的一个 $A$ 动机 $M$，其中 $A$ 为 $\mathbb Q$ 的有限扩张的有限直积。令 $\Lambda=A\otimes\mathbb Q_p$，以及
\[
 M_p^*(1)=\Hom_\Lambda(M_p,\Lambda)(1),
\]
其中 $\Gal(\overline{\mathbb Q}/\mathbb Q)$ 在 Hom 上按 $(\sigma,h)\mapsto h\circ\sigma^{-1}$ 作用。注意，由 3.2.5 中的信念，$M_p^*(1)$ 是 $A$ 动机 $\Hom_A(M,A)(1)$ 的 $p$ 进平展实现。令 $S$ 为包含 $p$ 的有限素数集，使 $M_p$ 在 $S$ 外不分歧。条件 3.2.2(v) 讨论 zeta 元素 $z_\Lambda(\Spec(\mathbb Z[1/S]),M_p)$ 与数值
\[
 \lim_{s\to0}s^{-e}L_{A,S}(M_p^*(1),s)
\]
之间的关系，其中 $e:\Spec(A\otimes\mathbb C)\to\mathbb Z$ 是 $L_S(M_p^*(1),s)$ 在 $s=0$ 处的阶。

3.2.2(v) 的完整描述见第 III 章 §2。这里讨论两个特殊情形，即\term{严格 K 理论情形}{strictly K-theoretic case}{cas strictement K-théorique}与\term{严格临界情形}{strictly critical case}{cas strictement critique}；从猜想的观点看，它们分别包含 §2.2 与 §2.3 中处理的情形。

令
\[
 \Sigma=M_B(-1)^+,\qquad H^i_S=H^i(\mathbb Z[1/S],M_p),
\]
其中 $(-1)$ 表示逆 Tate 扭转 $\otimes_{\mathbb Q}\mathbb Q(2\pi i)^{-1}$，$(\ )^+$ 表示 $\Gal(\mathbb C/\mathbb R)$ 不动部分。

在严格 K 理论情形与严格临界情形中，按定义均有 $H^i_S=0$（$i\ne1$），因此
\begin{equation*}\tag{1}
 \Delta_\Lambda(\Spec(\mathbb Z[1/S]),M_p)
 \simeq\det_\Lambda(H^1_S)\otimes_\Lambda\{\det_A(\Sigma)\}^{-1}.
\end{equation*}

\origpage{99}{94}
事实上，由 3.2.3(2)，有
\[
 \Sigma\otimes\mathbb Q_p=H^0(\mathbb R,M_p(-1))=R\Gamma(\mathbb R,M_p(-1)).
\]
例如，若 $A=\mathbb Q$，$M=H^m(X)(r)$，其中 $X$ 为 $\mathbb Q$ 上的光滑固有概形，且 $X,m,r,S$ 满足猜想 2.2.7（分别地，2.3.3）的假设，那么三元组
\[(\Spec(\mathbb Z[1/S]),\mathbb Q_p,M_p)\]
猜想应为严格 K 理论的（分别地，严格临界的）。对这一三元组，条件 3.2.2(v) 为
\[
 z_{\mathbb Q_p}(\Spec(\mathbb Z[1/S]),M_p)=\pm z_\gamma\otimes\gamma^{-1},
\]
其中 $\gamma,z_\gamma$ 如 2.2.7（分别地，2.3.3）。

先讨论严格临界情形。设 $M$ 为 $\mathbb Q$ 上的动机，令 $\Omega=M^0_{\dR}$。如下定义典范同态
\[
 a:\Omega\otimes\mathbb R\longrightarrow\Sigma\otimes\mathbb R,
 \qquad b:H^1_S\longrightarrow\Omega\otimes\mathbb Q_p.
\]
定义 $a$ 为复合映射
\begin{align*}
 \Omega\otimes\mathbb R&\hookrightarrow M_{\dR}\otimes\mathbb R
 \simeq(M_B\otimes\mathbb C)^+\\
 &\longrightarrow(M_B\otimes\mathbb C)^+/(M_B\otimes\mathbb R)^+
 \xrightarrow{\sim}(M_B(-1))^+\otimes\mathbb R,
\end{align*}
其中第一个同构来自 3.2.3(1)。\footnote{译注：此处及下句的引用原书如此；相应结构分别列于 3.2.4(4) 和 3.2.4(6)。}定义 $b$ 为复合映射
\[
 H^1_S\longrightarrow H^1(\mathbb Q_p,M_p)
 \longrightarrow D^0_{\dR}(M_p)\xrightarrow{\sim}M^0_{\dR}\otimes\mathbb Q_p,
\]
其中第二个箭头是第 II 章 §1.2 定义的“对偶指数映射”，同构来自 3.2.3(3)。

若同态 $a,b$ 都是双射，且对所有 $i\ne1$ 有 $H^i_S=0$，就称 $M$ 为严格临界的。我们猜想，若 $a,b$ 都是双射，则最后一个条件自动成立。

在严格临界情形中，条件 3.2.2(v) 变为下面的 (2)、(3)。

\noindent(2) $L_{A,S}(M_p^*(1),s)$ 在 $s=0$ 处的阶，即一个函数 $\Spec(A\otimes\mathbb C)\to\mathbb Z$，恒为零。

\origpage{100}{95}
\noindent(3) 考虑分别由 $a,b$ 诱导的同构
\begin{align*}
 \det_A(\Omega)\otimes_A\det_A(\Sigma)^{-1}\otimes\mathbb R
 &\xrightarrow{\sim}A\otimes\mathbb R,\\
 \det_A(\Omega)\otimes_A\det_A(\Sigma)^{-1}\otimes\mathbb Q_p
 &\xrightarrow{\sim}\Delta_\Lambda(\Spec(\mathbb Z[1/S]),M_p).
\end{align*}
那么 $z_\Lambda(\Spec(\mathbb Z[1/S]),M_p)$ 来自 $\det_A(\Omega)\otimes_A\det_A(\Sigma)^{-1}$ 中的一个元素，且该元素在 $A\otimes\mathbb R$ 中的像为 $\pm L_{A,S}(M_p^*(1),0)$。这里 $\pm$ 表示一个函数 $\Spec(A)\to\{\pm1\}$。

下面讨论严格 K 理论情形。为此须引入与 $M$ 关联的 K 群 $\Psi$。当 $M=H^m(X)(r)$，其中 $X$ 是 $\mathbb Q$ 上的光滑固有概形、$m,r\in\mathbb Z$ 时，$\Psi$ 按 2.2.4 定义。一般情形中，一个猜想性的构造是
\[
 \Psi=\Ext^1(\mathbb Q,M),
\]
其中 $\Ext^1$ 是猜想中的 $\mathbb Q$ 上动机的阿贝尔范畴中的扩张群。在我们的情形中，$M$ 为 $A$ 动机，因此 $\Psi$ 应带有 $A$ 模结构。此外，猜想应有一个调节子映射，当 $\Omega=0$ 时，它是一个 $A\otimes\mathbb R$ 同态
\[
 a:\Psi\otimes\mathbb R\longrightarrow\Sigma\otimes\mathbb R.
\]
还应有一个陈类映射，它是一个 $\Lambda$ 同态
\[
 b:\Psi\otimes\mathbb Q_p\longrightarrow H^1_S.
\]
若 $M=H^m(X)(r)$，这些 $a,b$ 就是 §2.1 中的映射。\footnote{译注：引用原书如此；相应映射的定义见 §2.2，特别是 2.2.5。}一般情形见 [FP$_1$] 或第 III 章 §2。若 $\Omega=0$，$a,b$ 都是同构，且对 $i\ne1$ 有 $H^i_S=0$，则称 $M$ 为严格 K 理论的。我们猜想，$\Omega=0$ 且 $a,b$ 为双射时，最后一个条件自动成立。

在严格 K 理论情形中，条件 3.2.2(v) 变为下面的 (4)、(5)。

\origpage{101}{96}
\noindent(4) $L_{A,S}(M_p^*(1),s)$ 在 $s=0$ 处的阶 $e$，作为函数 $\Spec(A\otimes\mathbb C)\to\mathbb Z$，等于复合映射
\[
 \Spec(A\otimes\mathbb C)\longrightarrow\Spec(A)\longrightarrow\mathbb Z,
\]
其中第二个箭头是 $\rank_A(\Sigma)$。

\noindent(5) 考虑分别由 $a,b$ 诱导的同构
\begin{align*}
 \det_A(\Psi)\otimes_A\det_A(\Sigma)^{-1}\otimes\mathbb R
 &\xrightarrow{\sim}A\otimes\mathbb R,\\
 \det_A(\Psi)\otimes_A\det_A(\Sigma)^{-1}\otimes\mathbb Q_p
 &\xrightarrow{\sim}\Delta_\Lambda(\Spec(\mathbb Z[1/S]),M_p).
\end{align*}
那么 $z_\Lambda(\Spec(\mathbb Z[1/S]),M_p)$ 来自 $\det_A(\Psi)\otimes_A\det_A(\Sigma)^{-1}$ 中的一个元素，且该元素在 $A\otimes\mathbb R$ 中的像为
\[
 \pm\lim_{s\to0}s^{-e}L_{A,S}(M_p^*(1),s).
\]
这里 $\pm$ 表示一个函数 $\Spec(A)\to\{\pm1\}$。

\subsection*{§3.3. 可逆模中 Galois 表示的 zeta 元素}
令 $p$ 为固定的奇素数。

\paragraph{3.3.1.} 本节考虑 3.2.1 中的三元组 $(X,\Lambda,\mathcal F)$，其中 $X=\Spec(\mathbb Z[1/S])$，$S$ 为包含 $p$ 的有限素数集，而 $\mathcal F$ 为可逆光滑 $\Lambda$ 层，作为 $D_{\mathrm{ctf}}(X,\Lambda)$ 的对象放在次数 $0$。

我们将在 $S=\{p\}$ 时定义 $\Delta_\Lambda(X,\mathcal F)$ 的典范 $\Lambda$ 基 $z_\Lambda(X,\mathcal F)$，并给出一般 $S$ 的定义方案。

经典岩泽理论与\term{分圆单位}{cyclotomic unit}{unité cyclotomique}在本节起核心作用。

本节不问这里讨论的 $z_\Lambda(X,\mathcal F)$ 是否满足 3.2.2(v) 的性质。这个问题将在后面讨论（见 3.4.4），并将得到部分肯定的回答。

本节的写作受到 Soulé [So$_2$] 与 Deligne [De$_5$] 关于“$p$ 进 Galois 上同调中的分圆元素”的定义的影响。

\origpage{102}{97}
\paragraph{3.3.2.} 我们定义 $z_\Lambda(X,\mathcal F)$ 的方案，是从下列“泛情形”推出它。

固定 $\mathbb Q$ 的一个代数闭包 $\overline{\mathbb Q}$，令 $\mathbb Q^{\mathrm{ab},S}$ 为 $\overline{\mathbb Q}$ 中所有在 $S$ 外不分歧的 $\mathbb Q$ 的有限阿贝尔扩张之并，令
\[
 G=\Gal(\mathbb Q^{\mathrm{ab},S}/\mathbb Q).
\]
所谓“泛情形”，是指三元组 $(X,\Lambda_0,\mathcal F_0)$，其中 $X=\Spec(\mathbb Z[1/S])$，且
\begin{equation*}\tag{1}
 \Lambda_0=\mathbb Z_p[[G]],\qquad
 \mathcal F_0=\varprojlim_L(f_L)_*f_L^*(\mathbb Z_p(1)_X).
\end{equation*}
这里 $L$ 遍历包含在 $\mathbb Q^{\mathrm{ab},S}$ 中的 $\mathbb Q$ 的所有有限扩张，$f_L:\Spec(\mathcal O_L[1/S])\to\Spec(\mathbb Z[1/S])$ 是典范态射，极限按迹映射取得。若有 $z_{\Lambda_0}(X,\mathcal F_0)$，则对任意如 3.1.1 的 $\Lambda$ 和任意 $X$ 上可逆光滑 $\Lambda$ 层 $\mathcal F$，可如下得到 $z_\Lambda(X,\mathcal F)$。

令 $\chi_{\mathcal F}:G\to\Lambda^\times$ 表示 $G$ 在 $H^0(\overline{\mathbb Q},\mathcal F)=H^0(\mathbb Q^{\mathrm{ab},S},\mathcal F)$ 上的作用。考虑同态
\[
 \Lambda_0\longrightarrow\Lambda,\qquad
 \sigma\longmapsto\chi_{\mathrm{cyclo}}(\sigma)\chi_{\mathcal F}(\sigma)^{-1}
 \quad(\sigma\in G).
\]
固定 $\Gamma(\overline{\mathbb Q},\mathcal F(-1))$ 的一个 $\Lambda$ 基 $\nu$，则在 $D_{\mathrm{ctf}}(X,\Lambda)$ 中有同构
\[
 h_\nu:\mathcal F_0\otimes_{\Lambda_0}\Lambda\xrightarrow{\sim}\mathcal F,
\]
定义如下。先设 $\Lambda$ 为有限环；一般 $\Lambda$ 的定义，可由此通过 $\varprojlim$ 和 $\otimes\mathbb Q$ 得到。此时存在 $L$，使 $f_L^*(\mathcal F)$ 为常值层。有同态
\[
 f_L^*(\mathbb Z_p(1)_X)\longrightarrow f_L^*(\mathcal F),\qquad x\longmapsto x\nu.
\]
它诱导映射
\[
 (f_L)_*f_L^*(\mathbb Z_p(1)_X)\longrightarrow(f_L)_*f_L^*(\mathcal F)
 \longrightarrow\mathcal F,
\]
\origpage{103}{98}
其中第二个箭头为迹映射。这一复合映射诱导所需同构 $h_\nu:\mathcal F_0\otimes_{\Lambda_0}\Lambda\xrightarrow{\sim}\mathcal F$。

假设有 $\Delta_{\Lambda_0}(X,\mathcal F_0)$ 的 $\Lambda_0$ 基 $z_{\Lambda_0}(X,\mathcal F_0)$。定义 $\Lambda$ 基 $z_\Lambda(X,\mathcal F)$ 为该元素在 $h_\nu$ 诱导的同构
\[
 \Delta_{\Lambda_0}(X,\mathcal F_0)\otimes_{\Lambda_0}\Lambda
 \xrightarrow{\sim}\Delta_\Lambda(X,\mathcal F)
\]
下的像。由 3.1.3(2)，此元素与 $\nu$ 的选择无关。事实上，若将 $\nu$ 换为 $a\nu$，其中 $a\in\Lambda^\times$，则 $z_\Lambda(X,\mathcal F)$ 乘以 $a^e$，其中
\[
 e=-\rank_\Lambda R\Gamma(X,\mathcal F)
 -\rank_\Lambda R\Gamma(\mathbb R,\mathcal F(-1))=0.
\]

\paragraph{3.3.3.} 本节余下部分，分别以 $\Lambda,\mathcal F$ 代替 3.3.2(1) 中的 $\Lambda_0,\mathcal F_0$。我们在 $S=\{p\}$ 时定义 $\Delta_\Lambda(X,\mathcal F)$ 的典范 $\Lambda$ 基 $z_\Lambda(X,\mathcal F)$，并给出一般 $S$ 的定义方案。注意
\[
 H^i(X,\mathcal F)=\varprojlim_N H^i(\mathbb Z[\alpha_N,1/S],\mathbb Z_p(1)),
\]
其中 $N$ 遍历所有满足“$S$ 恰为 $N$ 的素因子集”的整数 $N\ge1$，$\alpha_N$ 是 $\overline{\mathbb Q}$ 中的一个本原 $N$ 次单位根，极限按迹映射取得。当 $i\ne1,2$ 时，这个群为零。还应注意
\[
 H^i(\mathbb R,\mathcal F(-1))
 =\varprojlim_N H^i(\mathbb R\otimes\mathbb Q(\alpha_N),\mathbb Z_p),
\]
其中 $N,\alpha_N$ 如上。当 $i\ne0$ 时，这个群为零；$H^0(\mathbb R,\mathcal F(-1))$ 是秩为 $1$ 的自由 $\Lambda/(1-\tau)$ 模，其中 $\tau\in G$ 表示复共轭。

\paragraph{3.3.4.} 令 $\tau\in G$ 为复共轭，令
\[
 \Lambda_+=\Lambda/(1-\tau)\Lambda,\quad
 \Lambda_-=\Lambda/(1+\tau)\Lambda,
 \qquad
 \mathcal F_+=\mathcal F/(1-\tau)\mathcal F,\quad
 \mathcal F_-=\mathcal F/(1+\tau)\mathcal F.
\]
\origpage{104}{99}
那么
\[
 \Lambda\xrightarrow{\sim}\Lambda_+\times\Lambda_-,\qquad
 \mathcal F\xrightarrow{\sim}\mathcal F_+\times\mathcal F_-,
\]
\[
 \Delta_\Lambda(X,\mathcal F)\xrightarrow{\sim}
 \Delta_{\Lambda_+}(X,\mathcal F_+)\times\Delta_{\Lambda_-}(X,\mathcal F_-).
\]
因此应分别考虑 $\varepsilon=\pm$ 对应的 $z_{\Lambda_\varepsilon}(X,\mathcal F_\varepsilon)$，并将 $z_\Lambda(X,\mathcal F)$ 定义为
\[(z_{\Lambda_+}(X,\mathcal F_+),z_{\Lambda_-}(X,\mathcal F_-))\]
在最后一个同构下的逆像。

\paragraph{3.3.5.} 作为预备步骤，我们不假设任何猜想，如下定义元素
\[
 z_{\Lambda_+}(X,\mathcal F_+)
 \in H^1(\mathbb Z[1/S],\mathcal F_+)
 \otimes_{\Lambda_+}H^0(\mathbb R,\mathcal F_+(-1))^{-1},
 \qquad
 z_{\Lambda_-}(X,\mathcal F_-)\in Q(\Lambda_-).
\]
这里 $H^0(\mathbb R,\mathcal F_+(-1))^{-1}$ 为可逆 $\Lambda_+$ 模 $H^0(\mathbb R,\mathcal F_+(-1))$ 的逆模，$Q(\Lambda_-)$ 表示 $\Lambda_-$ 的\term{全商环}{total quotient ring}{anneau total des fractions}。

先考虑 $z_{\Lambda_+}(X,\mathcal F_+)$。对整数 $N\ge2$，令 $c_N$ 为如下定义的“分圆单位”：
\[
 c_N=(1-\alpha)(1-\alpha^{-1})\in\mathbb R^\times,
 \qquad\alpha=\exp(2\pi iN^{-1}).
\]
$c_N$ 是 $\mathbb Q(\alpha)^+$ 的整数环 $\mathbb Z[\alpha]^+$ 中的单位（因而是代数数论通常意义下的“单位”）当且仅当 $N$ 为合数\translatornote{9}；若 $N$ 为素数幂，它是 $\mathbb Z[\alpha,1/N]^+$ 中的单位。这里 $(\ )^+$ 表示 $\Gal(\mathbb C/\mathbb R)$ 不动部分。

定义 $z_{\Lambda_+}(X,\mathcal F_+)$ 为与下列 $\Lambda_+$ 模同态对应的元素：
\[
 H^0(\mathbb R,\mathcal F_+(-1))\longrightarrow H^1(\mathbb Z[1/S],\mathcal F_+).
\]
这一同态定义为以下映射的逆极限：
\begin{align*}
 H^0(\mathbb Q(\alpha_N)^+\otimes\mathbb R,\mathbb Z_p)
 &=\bigoplus_\iota H^0(\mathbb R,\mathbb Z_p)
 \longrightarrow H^1(\mathbb Z[\alpha_N,1/S]^+,\mathbb Z_p(1)),\\
 (a_\iota)_\iota&\longmapsto
 \sum_\iota2^{-1}a_\iota\{\iota^{-1}(c_N)\}.
\end{align*}

\origpage{105}{100}
这里 $a_\iota\in H^0(\mathbb R,\mathbb Z_p)$；$N$ 遍历所有满足“$S$ 恰为 $N$ 的素因子集”的整数 $N\ge1$，$\alpha_N$ 为一个本原 $N$ 次单位根。对每个 $N$，$\iota$ 遍历所有同态 $\mathbb Q(\alpha_N)^+\to\mathbb R$。对 $x\in(\mathbb Z[\alpha_N,1/S]^+)^\times$，$\{x\}$ 表示它通过 Kummer 正合列
\[
 0\longrightarrow\mathbb Z/p^i(1)\longrightarrow\mathbb G_m
 \xrightarrow{p^i}\mathbb G_m\longrightarrow0
\]
的连接映射得到的 $H^1(\mathbb Z[\alpha_N,1/S]^+,\mathbb Z_p(1))$ 中的像。

接着考虑 $z_-(X,\mathcal F_-)\in Q(\Lambda_-)$。令
\[
 z=\varprojlim_N\sum_{a\in(\mathbb Z/N)^\times}\zeta_{a(N)}(0)\sigma_a^{-1}
 \in\varprojlim_N\mathbb Q[\Gal(\mathbb Q(\alpha_N)/\mathbb Q)],
\]
其中 $N$ 遍历所有满足“$S$ 恰为 $N$ 的素因子集”的整数 $N\ge1$，而 $\sigma_a\in\Gal(\mathbb Q(\alpha_N)/\mathbb Q)$ 满足 $\sigma_a(\alpha_N)=\alpha_N^a$。这个 $z$ 可看作 $Q(\Lambda)$ 的元素，并且满足
\[
 (1-\chi_{\mathrm{cyclo}}(\sigma)\sigma^{-1})z\in\Lambda
 \qquad\bigl(\sigma\in\Gal(\mathbb Q^{\mathrm{ab},S}/\mathbb Q)\bigr)
\]
（1.3.7(1)）。定义 $z_-(X,\mathcal F_-)$ 为 $z\in Q(\Lambda)$ 在 $Q(\Lambda_-)$ 中的像。

\paragraph{3.3.6.} 这里指出，分圆单位是与 zeta 值密切相关的“伟大元素”，如下面的 (1) 所示；上面对 $z_+(X,\mathcal F_+)$ 与 $z_-(X,\mathcal F_-)$ 的定义也存在平行关系。对 $\alpha=\exp(2\pi iN^{-1})$（$N\ge2$）和 $\mathbb Z/N$ 中非零元素 $a$，有
\begin{equation*}\tag{1}
 -2^{-1}\log((1-\alpha^a)(1-\alpha^{-a}))
 =\lim_{s\to0}s^{-1}\bigl(\zeta_{a(N)}(s)+\zeta_{-a(N)}(s)\bigr).
\end{equation*}

\paragraph{3.3.7.} 现在尝试把 $z_{\Lambda_\varepsilon}(X,\mathcal F_\varepsilon)$（$\varepsilon=\pm$）看作 $\Delta_{\Lambda_\varepsilon}(X,\mathcal F_\varepsilon)$ 的 $\Lambda_\varepsilon$ 基。由于一般的 $\Lambda_\varepsilon$ 不是 Noether 环，而加上 Noether 假设会使某些论证更简单，我们考虑 Noether 环
\[
 \Lambda_{Np^\infty}=\mathbb Z_p[[G_{Np^\infty}]],
\]
\origpage{106}{101}
其中
\[
 G_{Np^\infty}=\varprojlim_n\Gal(\mathbb Q(\alpha_{Np^n})/\mathbb Q),
\]
$N\ge1$ 为整数，且 $S$ 恰为 $Np$ 的素因子集。此时 $Q(\Lambda_{Np^\infty})$ 是有限个域的直积。令
\begin{align*}
 \Lambda_{Np^\infty,+}&=\Lambda_{Np^\infty}/(1-\tau)\Lambda_{Np^\infty},\qquad
 \Lambda_{Np^\infty,-}=\Lambda_{Np^\infty}/(1+\tau)\Lambda_{Np^\infty},\\
 \mathcal F_{Np^\infty}&=\mathcal F\otimes_\Lambda\Lambda_{Np^\infty},\qquad
 \mathcal F_{Np^\infty,\varepsilon}=\mathcal F\otimes_\Lambda\Lambda_{Np^\infty,\varepsilon}
 \quad(\varepsilon=\pm).
\end{align*}
已知：

\noindent(1) 当 $i=1$ 时，
\[
 H^i(X,\mathcal F_{Np^\infty,+})
 \otimes_{\Lambda_{Np^\infty,+}}Q(\Lambda_{Np^\infty,+})
\]
是一个可逆 $Q(\Lambda_{Np^\infty,+})$ 模；当 $i\ne1$ 时，它为零。

\noindent(2) 对所有 $i$，
\[
 H^i(X,\mathcal F_{Np^\infty,-})
 \otimes_{\Lambda_{Np^\infty,-}}Q(\Lambda_{Np^\infty,-})=0.
\]
事实上，$i\ge3$ 时 $H^i=0$ 自动成立（[Ma]）；$H^0=0$ 容易看出；$H^2\otimes Q=0$ 可归结到理想类群的有限性；关于 $H^1\otimes Q$ 的断言则由 3.1.3(2) 得到。

由 (1)、(2)，得到同构
\begin{equation*}\tag{3}
 \begin{aligned}
 &\Delta_{\Lambda_{Np^\infty,+}}(X,\mathcal F_{Np^\infty,+})
 \otimes_{\Lambda_{Np^\infty,+}}Q(\Lambda_{Np^\infty,+})\\
 &\quad\simeq H^1(X,\mathcal F_{Np^\infty,+})
 \otimes_{\Lambda_{Np^\infty,+}}H^0(\mathbb R,\mathcal F_{Np^\infty,+}(-1))^{-1}
 \otimes_{\Lambda_{Np^\infty,+}}Q(\Lambda_{Np^\infty,+}),
 \end{aligned}
\end{equation*}
\begin{equation*}\tag{4}
 \Delta_{\Lambda_{Np^\infty,-}}(X,\mathcal F_{Np^\infty,-})
 \otimes_{\Lambda_{Np^\infty,-}}Q(\Lambda_{Np^\infty,-})
 \xrightarrow{\sim}Q(\Lambda_{Np^\infty,-}).
\end{equation*}
下面 3.3.8 的 $S=\{p\}$ 情形，就是 Mazur--Wiles [MW] 已证明的经典岩泽主猜想。

\begin{statement}{猜想 3.3.8}
令 $\varepsilon=+$（分别地，$-$）。那么 $z_\varepsilon(X,\mathcal F_\varepsilon)$（3.3.5）通过 3.3.7(3)（分别地，(4)）在
\[
 \Delta_{\Lambda_{Np^\infty,\varepsilon}}(X,\mathcal F_{Np^\infty,\varepsilon})
 \otimes_{\Lambda_{Np^\infty,\varepsilon}}Q(\Lambda_{Np^\infty,\varepsilon})
\]
中的像，是 $\Delta_{\Lambda_{Np^\infty,\varepsilon}}(X,\mathcal F_{Np^\infty,\varepsilon})$ 的一个 $\Lambda_{Np^\infty,\varepsilon}$ 基。
\end{statement}

\paragraph{3.3.9.} 事实上，当 $S=\{p\}$ 时，有 $\Lambda_\varepsilon=\Lambda_{Np^\infty,\varepsilon}$，且 $\Lambda_\varepsilon$ 为正则环。由 2.1.5，3.3.8 中 $\varepsilon=+$（分别地，$-$）的情形，等价于经典岩泽主猜想的熟悉形式
\origpage{107}{102}
\[
 \operatorname{char}_{\Lambda_+}(A/B)=\operatorname{char}_{\Lambda_+}(C),
\]
其中
\begin{align*}
 A&=H^1(X,\mathcal F_+)
 \simeq\varprojlim_n\bigl((\mathbb Z[\alpha_{p^n},1/p]^+)^\times\otimes\mathbb Z_p\bigr),\\
 B&=\im\bigl(H^0(\mathbb R,\mathcal F_+(-1))\longrightarrow H^1(X,\mathcal F_+)\bigr)
 \qquad\text{（3.3.5）},\\
 C&=H^2(X,\mathcal F_+)\simeq\varprojlim_n\Pic(\mathbb Z[\alpha_{p^n}]^+)\{p\}.
\end{align*}
两个同构均来自 Kummer 正合列。分别地，负部分的形式是
\[
 \Lambda_-\,z_{\Lambda_-}(X,\mathcal F_-)
 =\operatorname{char}_{\Lambda_-}(A)\operatorname{char}_{\Lambda_-}(B)^{-1},
\]
其中
\[
 A=H^2(X,\mathcal F_-)\simeq\varprojlim_n\Pic(\mathbb Z[\alpha_{p^n}])\{p\}^-,
 \qquad B=\operatorname{char}_{\Lambda_-}(\mathbb Z_p(1)).
\]
这里的同构来自 Kummer 正合列，上标 $-$ 表示复共轭按 $-1$ 作用的部分。\footnote{译注：原书如此；上式出现 $\operatorname{char}_{\Lambda_-}(B)$，而 $B$ 的定义又含一层 $\operatorname{char}_{\Lambda_-}$，此处保留两式原貌。}

这样，我们已在 $S=\{p\}$ 时定义 $\Delta_\Lambda(X,\mathcal F)$ 的 $\Lambda$ 基 $z_\Lambda(X,\mathcal F)$。对于一般的 $S$，当 $\Lambda_{Np^\infty,\varepsilon}$ 为正则环时，Mazur--Wiles [MW] 与 Wiles [Wi$_2$] 的结果可证明 3.3.8：由 2.1.5，3.3.8 可像上面那样改写为一个涉及“char”的等式，而这个等式已由 [MW][Wi$_2$] 证明。但若 $\Lambda_{Np^\infty,\varepsilon}$ 不是正则环，我还不能从 [MW][Wi$_2$] 的结果推出 3.3.8。

\subsection*{§3.4. 问题}
令 $p$ 为奇素数。

\paragraph{3.4.1.} 下面讨论由我们的岩泽主猜想自然产生的几个问题。

(1) §3.3 在某个特殊情形中给出的 zeta 元素的定义，是否与猜想 3.2.2 的条件 (v) 相容？参见 3.4.4。

\origpage{108}{103}
(2) 如何构造 zeta 元素？在这里我承认，我并不知道一个好的方法。参见 3.4.11。

(3) zeta 元素对算术代数几何有多大用处？参见 3.4.13。

\paragraph{3.4.2.} 在下面的 3.4.3--3.4.9 中，设 $(X,\Lambda,\mathcal F)$ 如 3.3.1，令 $A$ 为 $\mathbb Q$ 的有限扩张的有限直积，满足 $\Lambda=A\otimes\mathbb Q_p$。固定 $\mathbb Q$ 的代数闭包 $\overline{\mathbb Q}$，令 $\mathbb Q^{\mathrm{ab},S}$ 如 3.3.2。令
\[
 V=\Gamma(\Spec(\overline{\mathbb Q}),\mathcal F).
\]

\paragraph{3.4.3.} 动机的某些基本思想告诉我们：若 $\mathcal F$ 来自 $A$ 动机 $M$，则存在 $\mathbb Q^{\mathrm{ab},S}/\mathbb Q$ 的有限子扩张 $L$、同态
\[
 \eta:A'\overset{\mathrm{def}}=\mathbb Q[\Gal(L/\mathbb Q)]\longrightarrow A
\]
以及整数 $r\in\mathbb Z$，使得作为 $A$ 动机有
\[
 M\simeq H^0(\Spec(L))(r)\otimes_{A'}A.
\]
在这种情形中，
\[
 \chi_{\mathcal F}(\sigma)=\chi_{\mathrm{cyclo}}(\sigma)^r\eta(\sigma)^{-1}
 \qquad\bigl(\sigma\in\Gal(\mathbb Q^{\mathrm{ab},S}/\mathbb Q)\bigr),
\]
这里也用 $\eta$ 表示复合同态
\[
 \Gal(\mathbb Q^{\mathrm{ab},S}/\mathbb Q)\longrightarrow\Gal(L/\mathbb Q)
 \xrightarrow{\eta}A^\times.
\]

\paragraph{3.4.4.} 在 3.4.5--3.4.9 中，设 $r\in\mathbb Z$，并假设
\[
 \chi_{\mathcal F}(\sigma)=\chi_{\mathrm{cyclo}}(\sigma)^r\lambda(\sigma),
\]
其中同态 $\lambda:\Gal(\mathbb Q^{\mathrm{ab},S}/\mathbb Q)\to A^\times$ 相对于 $A^\times$ 的离散拓扑连续。

我们按如下方式讨论问题 3.4.1(1)。这里 $\tau\in\Gal(\mathbb Q^{\mathrm{ab},S}/\mathbb Q)$ 为复共轭。
\begin{center}
\begin{tabular}{ll}
$r\ge1$ 且 $\chi_{\mathcal F}(\tau)=1$ & 见 3.4.8；\\
$r\ge1$ 且 $\chi_{\mathcal F}(\tau)=-1$ & 见 3.4.9；\\
$r\le0$ 且 $\chi_{\mathcal F}(\tau)=1$ & 见第 III 章 1.2.6；\\
$r\le0$ 且 $\chi_{\mathcal F}(\tau)=-1$ & 见第 III 章 §3。
\end{tabular}
\end{center}

\origpage{109}{104}
考虑 2.2.3 的 $L$ 函数。容易看出 2.2.3 中的条件 (i)、(ii) 成立，并且有：

\begin{statement}{引理 3.4.5}
取整数 $N\ge1$，使 $S$ 恰为 $N$ 的素因子集，且 $\lambda$ 通过 $\Gal(\mathbb Q(\alpha_N)/\mathbb Q)$ 分解，其中 $\alpha_N$ 是 $\overline{\mathbb Q}$ 中的一个本原 $N$ 次单位根。那么
\[
 L_{A,S}(V^*(1),s)
 =\sum_{a\in(\mathbb Z/N)^\times}\zeta_{a(N)}(s+1-r)\lambda(\sigma_a),
\]
其中 $\sigma_a\in\Gal(\mathbb Q^{\mathrm{ab},S}/\mathbb Q)$ 满足 $\sigma_a(\alpha_N)=\alpha_N^a$。
\end{statement}

先考虑 $r\ge1$ 且 $\chi_{\mathcal F}(\tau)=1$ 的情形。由定理 1.3.7 与 3.4.5，有：

\begin{statement}{命题 3.4.6}
假设 $r\ge1$ 且 $\chi_{\mathcal F}(\tau)=1$。取任意 $\mu\in\Gal(\mathbb Q^{\mathrm{ab},S}/\mathbb Q)$，使 $\chi_{\mathrm{cyclo}}(\mu)$ 为无限阶。令 $z$ 为
\[
 \mathbb Z_p[[\Gal(\mathbb Q^{\mathrm{ab},S}/\mathbb Q)]]
 [(1-\chi_{\mathrm{cyclo}}(\mu)\mu^{-1})^{-1}]/(\tau+1)
\]
中在 3.3.5 记为 $z_-(X,\mathcal F_-)$ 的元素。那么，$z$ 在映射
\[
 \begin{gathered}
 \mathbb Z_p[[\Gal(\mathbb Q^{\mathrm{ab},S}/\mathbb Q)]]
 [(1-\chi_{\mathrm{cyclo}}(\mu)\mu^{-1})^{-1}]/(\tau+1)
 \longrightarrow\Lambda,\\
 \sigma\longmapsto\chi_{\mathrm{cyclo}}(\sigma)\chi_{\mathcal F}(\sigma)^{-1}
 \quad\bigl(\sigma\in\Gal(\mathbb Q^{\mathrm{ab},S}/\mathbb Q)\bigr)
 \end{gathered}
\]
下的像等于 $L_{A,S}(V^*(1),0)$。
\end{statement}

由 Soulé [So$_2$]，有：
\begin{statement}{命题 3.4.7}
假设 $r\ge2$ 且 $\chi_{\mathcal F}(\tau)=1$。那么对所有 $i\in\mathbb Z$，都有 $H^i(\mathbb Z[1/S],\mathcal F)=0$。
\end{statement}
事实上，对如 3.4.5 的 $N$，存在满射
\[
 \begin{gathered}
 H^i(\mathbb Z[\alpha_N,1/S],\mathbb Q_p(r))
 \simeq H^i(\mathbb Z[\alpha_N,1/S],\mathcal F)\\
 \longrightarrow H^i(\mathbb Z[1/S],\mathcal F),
 \end{gathered}
\]
最后一箭头为迹映射。由 Soulé [So$_1$]，$H^2(\mathbb Z[\alpha_N,1/S],\mathbb Q_p(r))=0$，故对所有 $i\ne1$，有 $H^i(\mathbb Z[1/S],\mathcal F)=0$；再由 3.1.3(2) 中的等式 (*)，得到 $H^1(\mathbb Z[1/S],\mathcal F)=0$。

\paragraph{3.4.8.} 假设 $r\ge2$、$\chi_{\mathcal F}(\tau)=1$，并假设 3.4.3 所述动机思想正确。于是，对于给出 $\mathcal F$ 的 $A$ 动机，有 $\Sigma=0$、$\Omega=0$。
\origpage{110}{105}
结合 3.4.7，可知此时处于严格临界情形（3.2.6）。

因此，3.4.6 表明：若猜想 3.3.8 成立，则典范同构 $\Delta_\Lambda(X,\mathcal F)\simeq\Lambda$ 把 $z_\Lambda(X,\mathcal F)$ 映为 $L_{A,S}(V^*(1),0)$；这一同构来自 $\det_\Lambda(0)\simeq\Lambda$。这就说明，在 $r\ge2$ 且 $\chi_{\mathcal F}(\tau)=1$ 的情形中，§3.3 的构造与 3.2.2(v) 相容。

在 $r=1$ 且 $\chi_{\mathcal F}(\tau)=1$ 的情形中，经过更细致的研究也能得到类似结论，但这里不作讨论。

\paragraph{3.4.9.} 下面考虑 $r\ge1$ 且 $\chi_{\mathcal F}(\tau)=-1$ 的情形。

假设 3.4.3 所述动机思想正确，考虑给出 $\mathcal F$ 的 $A$ 动机 $M$。于是 $M=M'\otimes_{A'}A$，其中
\[
 M'=H^0(\Spec(L))(r),\qquad A'=\mathbb Q[\Gal(L/\mathbb Q)],
\]
$L$ 是 $\mathbb Q^{\mathrm{ab},S}/\mathbb Q$ 的有限子扩张，并给定同态 $A'\to A$。有
\[
 M'_B=\bigoplus_\iota\mathbb Q(2\pi i)^r,\qquad M_B=M'_B\otimes_{A'}A,
\]
其中 $\iota$ 遍历所有嵌入 $L\to\mathbb C$。还应有
\[
 \Psi=\Psi'\otimes_{A'}A,\qquad\Psi'=K_{2r-1}(L)\otimes\mathbb Q.
\]
令
\[
 a:\Psi\otimes\mathbb R\longrightarrow\Sigma\otimes\mathbb R,
 \qquad\Sigma=M_B(-1)^+,
 \qquad b:\Psi\otimes\mathbb Q_p\longrightarrow H^1_{\mathrm{lim}}
\]
为分别由调节子映射
\[
 K_{2r-1}(L)\longrightarrow M'_B(-1)^+\otimes\mathbb R
\]
和陈类映射
\[
 K_{2r-1}(L)\longrightarrow\varinjlim_S H^1(\mathcal O_L[1/S],\mathbb Q_p(r))
\]
诱导的同态。

现在设 $r\ge2$。由 Borel [Bo] 与 Soulé [So$_2$]，$a,b$ 为双射，且 $H^2(\mathbb Z[1/S],\mathcal F)=0$，故处于严格 K 理论情形。此外，由 Borel [Bo]，$\Psi$ 是可逆 $A$ 模。
